\chapter{Presentación de análisis, resultados y discusión}
% Redacción con números finales (2026-08-10): 2x2 modelo x dominio, MAIRA-2, token de apertura,
% calibración y demos. Estilo impersonal UAX. Cada figura se justifica explícitamente en el texto.

En este capítulo se describe la muestra analizada, se presentan los resultados y se discute su
significado. La línea argumental es la siguiente: la atención del modelo está clínicamente anclada,
pero su morfología no permite estimar la fiabilidad; esta reside en la \emph{incertidumbre de la
decisión} del propio modelo y se detecta, sobre todo, para un tipo de error —las alucinaciones—. Se
demuestra que la señal es \emph{agnóstica al modelo} y \emph{robusta al dominio} validándola sobre
dos generadores (LLaVA-Rad y MAIRA-2) y dos dominios (Estados Unidos y España), y se traduce en una
confianza calibrada apta para el cribado clínico.

\section{Análisis descriptivo de la muestra}
El estudio de fiabilidad se articula como un diseño \(2\times2\) que cruza \textbf{dos generadores}
—LLaVA-Rad (7B, entrenado sobre MIMIC-CXR) y MAIRA-2 (7B, entrenado sobre MIMIC-CXR, PadChest y
USMix)— con \textbf{dos dominios}: Open-I (Indiana University, EE.\,UU.) y PadChest-GR (España). Para
cada celda se generan informes y se etiqueta cada uno como \emph{correcto}, \emph{alucinación}
(afirmar un hallazgo agudo ausente) u \emph{omisión} o \emph{miss} (perder un hallazgo agudo
presente). El etiquetado combina tres fuentes de verdad independientes para evitar la circularidad
de un único juez: diagnóstico estructurado oficial (MeSH de Open-I; hallazgos presentes/ausentes de
PadChest-GR) y juicio basado en modelo de lenguaje sobre la narrativa completa (\emph{G-Eval}). El
anclaje de la atención se evalúa aparte sobre las \textbf{984 anotaciones con caja} de NIH
ChestX-ray8. La Tabla~\ref{tab:muestra} resume las cuatro celdas. La prevalencia de error y el
predominio de informes normales son característicos del dominio; la cardiomegalia sobre-diagnosticada
en tórax normal aparece como sesgo sistemático de \emph{ambos} generadores.

\begin{table}[t]\centering
\caption{Muestra del diseño \(2\times2\) (etiquetas objetivas/G-Eval). n = informes; alu = alucinaciones; mis = omisiones.}
\label{tab:muestra}
\begin{tabular}{llrrr}
\toprule
Generador & Dominio & n & alu & mis \\ \midrule
LLaVA-Rad & Open-I (in-distribution)  & 800 & 110 & 228 \\
LLaVA-Rad & PadChest (OOD)            & 381 & 33  & 130 \\
MAIRA-2   & PadChest (in-distribution)& 379 & 27  & 80  \\
MAIRA-2   & Open-I (OOD)              & 300 & 12  & 22  \\ \bottomrule
\end{tabular}
\end{table}

\section{El generador propio y la lección sobre las métricas}
El generador propio se desarrolló de forma deliberada como paso previo a emplear un modelo del
estado del arte ya entrenado: construirlo desde cero permite observar su \textbf{funcionamiento
interno} —cómo se integra la señal visual en la generación de texto, y dónde falla— de un modo que
un modelo de caja cerrada no ofrece, y expone de primera mano las \textbf{limitaciones propias de un
modelo de menor capacidad} entrenado con datos escasos. En su implementación se documentaron y
resolvieron dos fallos característicos. El primero, el \textbf{colapso de la generación} al inyectar la atención
cruzada: la contribución visual, inicialmente ruidosa, desestabilizaba el decodificador preentrenado;
se mitigó con la \textbf{compuerta inicializada a cero}, que preserva el comportamiento del modelo de
lenguaje y deja que la señal visual se incorpore de forma gradual. El segundo, la \textbf{ceguera a la
imagen}: el decodificador tendía a ignorar la imagen y a reproducir la plantilla mayoritaria; se
mitigó congelando el decodificador base y entrenando solo la atención cruzada (PEFT).

Sobre esta base se diseñó una ablación (Tabla~\ref{tab:ablacion}) que introduce progresivamente la
destilación de conocimiento textual y visual, con el objetivo de medir su efecto sobre la calidad y el
anclaje. La lección metodológica más relevante emergió al evaluar estas variantes: la métrica
\textbf{ROUGE engaña}. Una variante mejor anclada a la imagen —que describe hallazgos reales— puede
obtener un ROUGE \emph{inferior} al de un modelo ciego que memoriza la plantilla ``normal'', porque en
un dominio dominado por informes normales el solapamiento léxico premia la reproducción de la
plantilla, no la lectura de la imagen. Este hallazgo, coherente con la crítica del Capítulo~3, motivó
el abandono de las métricas de solapamiento en favor de la \textbf{corrección clínica} y, en última
instancia, el giro de todo el trabajo hacia la \textbf{fiabilidad}. Junto con las limitaciones de
capacidad ya observadas en el propio modelo, esta lección sobre las métricas es la que llevó a
\textbf{elevar el análisis a un MLLM del estado del arte} (LLaVA-Rad): solo un generador de calidad
suficiente permite separar los errores reales de los artefactos de un modelo de escala reducida, y
solo evaluando la corrección clínica —no el solapamiento léxico— tiene sentido preguntarse si la
explicación o la incertidumbre predicen la fiabilidad.

\begin{table}[t]\centering
\caption{Ablación de destilación del generador propio (M0--M3): fuentes de conocimiento introducidas.}
\label{tab:ablacion}
\begin{tabular}{lll}
\toprule
Variante & Destilación textual & Destilación visual \\ \midrule
M0 (línea base) & --- & --- \\
M1 & CXR-BERT (torre congelada) & --- \\
M2 & CXR-BERT (destilación al \emph{student}) & --- \\
M3 & CXR-BERT (destilación) & RAD-DINO (CLS $\to$ CLS) \\ \bottomrule
\end{tabular}
\end{table}

\section{La atención está clínicamente anclada}
Contrastada con las cajas reales de NIH, la atención de LLaVA-Rad se concentra sobre la lesión muy
por encima del azar: \textbf{lift global 3,03} y \emph{pointing-game} 0,38 (3,7 veces el azar). El
anclaje varía con sentido clínico (Tabla~\ref{tab:nih}): fuerte en hallazgos focales pequeños
(nódulo, masa), débil en la cardiomegalia, difusa y de caja grande. Este resultado justifica el uso
del mapa de calor como puntero de \emph{dónde} mira el modelo, pero —como se verá— no como medida de
\emph{cuánto} fiarse.

\begin{table}[t]\centering
\caption{Anclaje objetivo sobre NIH (n = 388 hallazgos detectados). Lift = energía-en-caja / azar.}
\label{tab:nih}
\begin{tabular}{lrrr}
\toprule
Patología & Detección & Lift & Pointing \\ \midrule
Nódulo & 0,22 & 7,91 & 0,47 \\ Masa & 0,18 & 5,10 & 0,87 \\
Atelectasia & 0,29 & 3,45 & 0,30 \\ Derrame & 0,58 & 3,23 & 0,36 \\
Neumotórax & 0,74 & 2,81 & 0,37 \\ Neumonía & 0,29 & 2,67 & 0,23 \\
Infiltrado & 0,19 & 2,62 & 0,48 \\ Cardiomegalia & 0,58 & 1,63 & 0,39 \\ \midrule
\textbf{Global} & 0,39 & \textbf{3,03} & 0,38 \\ \bottomrule
\end{tabular}
\end{table}

La Figura~\ref{fig:nih} ilustra visualmente este anclaje: para cada patología se muestran los casos de
mayor \emph{lift}, con el mapa de calor superpuesto y la caja real. Se observa que la atención se
concentra sobre la lesión, de forma especialmente nítida en los hallazgos focales pequeños (nódulo,
masa, neumotórax) y más difusa en la cardiomegalia, de caja grande.

\begin{figure}[p]\centering
\includegraphics[height=0.92\textheight]{fig_nih_montage.png}
\caption{Anclaje de la atención sobre NIH por patología (casos de mayor \emph{lift}). Mapa de calor
superpuesto y caja real (verde); se indica el \emph{lift}, el acierto/fallo del veredicto y el área
relativa de la caja. La atención se concentra sobre la lesión, sobre todo en hallazgos focales.}
\label{fig:nih}
\end{figure}

A modo de contraste, la Figura~\ref{fig:nihmed} muestra los casos de \emph{lift} mediano por patología:
el anclaje sigue siendo apreciable pero más difuso, ilustrando que la calidad de la localización varía
de forma continua y que no todos los aciertos van acompañados de una atención igualmente concentrada.

\begin{figure}[p]\centering
\includegraphics[height=0.92\textheight]{fig_nih_median.png}
\caption{Anclaje sobre NIH por patología en casos de \emph{lift} mediano. La atención sigue sobre la
región de la lesión pero de forma más difusa que en los casos de la Figura~\ref{fig:nih}.}
\label{fig:nihmed}
\end{figure}

\section{La morfología de la explicación no predice la fiabilidad}
Ninguna característica geométrica del mapa de calor (concentración, dispersión, agrupamiento) separa
de forma robusta los informes correctos de los erróneos (AUC 0,54; no significativo). Las señales
que aparentaban serlo resultaron artefactos detectados al validarlas: la aparente relación entre
anclaje (\emph{lift}) y acierto \textbf{desaparece} (AUC 0,52) al emplear un mapa de calor homogéneo,
revelándose como un artefacto de contrastar mapas por-hallazgo con mapas globales. El resultado se
confirma por tercera vía con las \textbf{cajas de anclaje propias de MAIRA-2}: el modelo localiza sus
hallazgos \emph{inventados} con cajas anatómicamente plausibles (una caja sobre el corazón para una
cardiomegalia inexistente, sobre el pulmón para una enfermedad intersticial inexistente), de modo que
la presencia y la forma del anclaje acompañan a la \emph{afirmación}, no al \emph{acierto} (frecuencia
de caja del 19\,\% en alucinaciones frente al 7\,\% en informes correctos, dirección inútil como
señal). En resumen, ni el mapa de calor de atención, ni el \emph{lift} sobre NIH, ni las cajas
entrenadas del propio modelo permiten leer la fiabilidad en la morfología de la explicación.

\section{La incertidumbre de la decisión detecta las alucinaciones}
Frente a la morfología, la \textbf{incertidumbre generativa} —operacionalizada como la probabilidad
mínima de \emph{contenido} de los tokens del informe— sí discrimina la alucinación del acierto. La
Figura~\ref{fig:sep} muestra, para las cuatro celdas, las distribuciones de esta señal en informes
correctos, con omisión y con alucinación. \emph{Se incluye esta figura porque hace visible el
resultado central}: las alucinaciones (rojo) se desplazan hacia la baja confianza y se separan de los
informes correctos (verde), con la única excepción de LLaVA-Rad fuera de distribución, donde ambas
distribuciones se solapan.

\begin{figure}[t]\centering
\includegraphics[width=\linewidth]{fig_separacion_clases.png}
\caption{Separación de clases de la probabilidad mínima de contenido (menor $=$ más incierto) en las
cuatro celdas del diseño. Verde: correcto; gris: omisión; rojo: alucinación; líneas discontinuas:
medias. Las alucinaciones se concentran en la zona incierta salvo en LLaVA-Rad$\cdot$PadChest (OOD).}
\label{fig:sep}
\end{figure}

\paragraph{El token de apertura contamina la señal (aportación metodológica).}
La probabilidad mínima calculada sobre \emph{todo} el informe rinde de forma mediocre porque el
mínimo recae casi siempre en el token de \emph{apertura} del informe (``The'', ``Frontal''\dots),
que refleja incertidumbre de estilo, no de contenido, y es uniformemente baja en todos los casos.
Excluir ese primer token —usar la probabilidad mínima de contenido— eleva sustancialmente la
detección en las cuatro celdas (Figura~\ref{fig:token}). \emph{Se justifica esta figura porque aísla
una decisión de diseño con impacto grande y generalizable}: el AUC pasa de 0,71 a 0,90 en Open-I y de
0,50 a 0,80 en MAIRA-2$\cdot$PadChest, sin coste alguno. La interpretación de este efecto es
instructiva: el primer token del informe concentra la incertidumbre \emph{estilística} —qué fórmula de
apertura emplear (``The'', ``Frontal''\dots)—, que es alta y uniforme en todos los informes con
independencia de su contenido, y que por tanto \emph{no discrimina} y solo añade ruido al mínimo. Al
restringir la señal al contenido, esta pasa a reflejar la incertidumbre \emph{diagnóstica} del modelo.
Más allá de la mejora numérica, el hallazgo ilustra un principio general aplicable a otras tareas de
generación: al resumir la incertidumbre de una secuencia conviene separar la duda sobre la forma de la
duda sobre el fondo, pues solo la segunda informa sobre la fiabilidad del contenido.

\begin{figure}[t]\centering
\includegraphics[width=0.92\linewidth]{fig_token_apertura.png}
\caption{Efecto de excluir el token de apertura. Barras grises: probabilidad mínima sobre todo el
informe; barras verdes: sobre el contenido (sin el primer token). La mejora es sistemática en las
cuatro celdas.}
\label{fig:token}
\end{figure}

\paragraph{La señal es agnóstica al modelo y robusta al dominio.}
La Figura~\ref{fig:2x2} sintetiza el resultado principal como matriz \(2\times2\). \emph{Se incluye
porque resume, en una sola imagen, la validez externa del método}: la detección es fuerte in-distribution
(LLaVA-Rad 0,90; MAIRA-2 0,80) y, crucialmente, \textbf{se mantiene fuera de distribución para el
modelo de mayor calidad} (MAIRA-2 0,81, prácticamente igual que su valor in-distribution). La única
celda débil es LLaVA-Rad fuera de distribución (0,61): un generador limitado, al enfrentarse a
imágenes ajenas a su entrenamiento, recae en un informe-plantilla emitido con confianza, lo que
erosiona la señal. La degradación fuera de distribución es, por tanto, una limitación del generador
débil, no del método. Este matiz es importante porque invierte la conclusión que se habría extraído de
haber estudiado un único generador: con LLaVA-Rad aislado, la caída a 0,61 fuera de distribución
sugeriría que la fiabilidad por incertidumbre no generaliza; el diseño cruzado revela que sí lo hace
cuando el modelo es lo bastante bueno, y que el problema reside en que un generador débil, ante
imágenes ajenas a su entrenamiento, deja de leer la imagen y ``rellena'' con su plantilla más
probable, emitida con confianza. La lección práctica es doble: conviene, primero, elegir un generador
de calidad, y segundo, calibrar y verificar la señal en cada dominio de despliegue antes de confiar en
ella.

\begin{figure}[t]\centering
\includegraphics[width=0.72\linewidth]{fig_2x2_auc.png}
\caption{AUC de detección de alucinaciones (probabilidad mínima de contenido) en el diseño
\(2\times2\). Verde $=$ alto. La señal es robusta salvo en el generador débil fuera de distribución.}
\label{fig:2x2}
\end{figure}

Las curvas ROC de la Figura~\ref{fig:roc} caracterizan el compromiso sensibilidad--especificidad de
la señal en las cuatro celdas. \emph{Se muestran porque traducen el AUC agregado en el comportamiento
operativo}: las tres celdas robustas dominan claramente la diagonal de azar, mientras que
LLaVA-Rad$\cdot$PadChest se aproxima a ella, coherente con el solapamiento de la Figura~\ref{fig:sep}.

\begin{figure}[t]\centering
\includegraphics[width=0.62\linewidth]{fig_roc.png}
\caption{Curvas ROC de detección de alucinaciones (probabilidad mínima de contenido). Tres de las
cuatro celdas dominan la diagonal de azar; LLaVA-Rad$\cdot$PadChest (OOD) se aproxima a ella.}
\label{fig:roc}
\end{figure}

\subsection{Ninguna señal adicional generaliza: en favor de la simplicidad}
Se exploró si combinar la probabilidad mínima de contenido con otras señales agnósticas (entropía,
verosimilitud media, auto-consistencia, dependencia de imagen) mejoraba la detección. La exploración
se realizó exigiendo que cualquier mejora \emph{generalizara} sobre los tres conjuntos con suficiente
número de alucinaciones. La Tabla~\ref{tab:senales} muestra el AUC univariante de cada señal por
conjunto: ninguna supera de forma consistente a la probabilidad mínima de contenido. Señales fuertes
en los conjuntos de LLaVA-Rad (entropía, verosimilitud) resultan inútiles en MAIRA-2, y la dependencia
de imagen ayuda en MAIRA-2 pero se invierte en Open-I; son, por tanto, específicas de modelo o de
dominio. En consecuencia, la fórmula final privilegia la \textbf{simplicidad robusta} de una única
señal frente a un combinado sobreajustado que no transferiría.

\begin{table}[t]\centering
\caption{AUC univariante de cada señal (orientada: mayor $=$ más probable alucinación) por conjunto.
Ninguna mejora de forma consistente a la probabilidad mínima de contenido.}
\label{tab:senales}
\begin{tabular}{lrrr}
\toprule
Señal & Open-I (LLaVA) & PadChest (LLaVA) & PadChest (MAIRA) \\ \midrule
\textbf{min\_prob de contenido} & \textbf{0,90} & \textbf{0,57} & \textbf{0,80} \\
Percentil 5 de probabilidad & 0,90 & 0,61 & 0,58 \\
Verosimilitud media & 0,89 & 0,65 & 0,49 \\
Entropía media & 0,92 & 0,65 & 0,46 \\
Auto-consistencia & 0,73 & 0,43 & 0,65 \\
Dependencia de imagen (blank) & 0,32 & 0,51 & 0,70 \\ \bottomrule
\end{tabular}
\end{table}

\section{Disociación alucinación--omisión}
La detección es \emph{asimétrica}. La Figura~\ref{fig:diso} contrasta, por celda, el AUC para separar
alucinaciones frente a correctos y omisiones frente a correctos. \emph{Se justifica esta figura porque
delimita con honestidad el alcance del método}: la señal detecta las alucinaciones de forma robusta
(0,80--0,90), pero las omisiones solo de forma parcial y variable (0,47--0,85). Las omisiones no son
siempre errores ``confiados''; como se aprecia en la Figura~\ref{fig:sep} (\mbox{distribución} gris), tienden
a situarse \emph{entre} los correctos y las alucinaciones, cerca del umbral, y fuera de distribución
(MAIRA-2$\cdot$Open-I) se desplazan hacia la incertidumbre. La interpretación es que, cuando el modelo
fabrica un hallazgo, titubea; cuando lo omite emitiendo un informe normal, a menudo lo hace con
confianza próxima a la de un informe normal correcto, aunque no siempre. En consecuencia, la fórmula
es un \emph{detector de alucinaciones} que captura además una fracción de omisiones como beneficio
colateral, no un detector general de error.

\begin{figure}[t]\centering
\includegraphics[width=0.82\linewidth]{fig_disociacion.png}
\caption{Disociación: AUC para separar alucinaciones (rojo) y omisiones (gris) de los informes
correctos, por celda. La incertidumbre detecta alucinaciones de forma robusta y omisiones solo
parcialmente.}
\label{fig:diso}
\end{figure}

\section{Calibración y confianza operativa}
Para convertir la señal en una salida clínica se estudia su calibración. La Figura~\ref{fig:cal}
representa la tasa de alucinación observada en función de la probabilidad mínima de contenido:
\emph{se incluye porque es el fundamento empírico del ``semáforo'' de confianza}, al mostrar una
relación monótona —a menor confianza del modelo, mayor tasa de alucinación observada— sobre la que se
definen bandas de riesgo. Sobre Open-I, las bandas por cuartiles del grupo correcto arrojan una tasa
de alucinación del 30\,\% en la banda roja (revisar), 6\,\% en la amarilla y $\approx$0\,\% en las
verdes. Como priorizador de revisión (Tabla~\ref{tab:oper}), un umbral de Youden captura el 84\,\% de
las alucinaciones —y, además, el 44\,\% de las omisiones, como beneficio colateral— revisando en
torno a un tercio de los estudios; a especificidad del 95\,\%, el 53\,\% de las alucinaciones revisando solo el 5\,\%. La
propuesta no es, pues, un clasificador binario —la baja prevalencia lo haría poco preciso— sino una
\textbf{confianza calibrada con intervalo} que ordena la carga de revisión, con el radiólogo como
decisor.

\begin{figure}[t]\centering
\includegraphics[width=0.66\linewidth]{fig_calibracion.png}
\caption{Calibración: tasa de alucinación observada frente a la probabilidad mínima de contenido
(media por decil, datos agregados). La relación monótona fundamenta las bandas del semáforo de
confianza.}
\label{fig:cal}
\end{figure}

\begin{table}[t]\centering
\caption{Puntos de operación sobre Open-I (verdad MeSH; puntuación de contenido). FA $=$ falsas alarmas.}
\label{tab:oper}
\begin{tabular}{lrrr}
\toprule
Umbral & Recall alucinación & Recall omisión & FA \\ \midrule
Youden          & 84\,\% & 44\,\% & 15\,\% \\
Especificidad 90\,\% & 71\,\% & 33\,\% & 10\,\% \\
Especificidad 95\,\% & 53\,\% & 20\,\% & 5\,\%  \\ \bottomrule
\end{tabular}
\end{table}

\section{Traducción clínica: la confianza como semáforo}
Los resultados anteriores solo tienen valor si pueden \emph{presentarse} al clínico de una forma que
mejore su trabajo sin introducir nuevos riesgos. Esta sección justifica las decisiones de diseño de la
demostración desarrollada y su encaje en un entorno médico real.

\subsection{Por qué una demostración interactiva}
Se ha optado por materializar el sistema en una \textbf{demostración interactiva} (con la biblioteca
Gradio) que, ante una radiografía, muestra en un único panel tres elementos: el \textbf{informe}
generado, el \textbf{anclaje} (dónde mira el modelo, mediante mapa de calor o cajas) y la
\textbf{confianza} calibrada (cuánto fiarse). La elección responde a tres criterios. Primero, la
\textbf{reproducibilidad y accesibilidad}: una interfaz web ligera permite que un radiólogo,
un investigador o un evaluador reproduzcan el comportamiento del sistema sin conocimientos técnicos ni
instalación compleja. Segundo, la \textbf{integración de las tres señales} en una sola vista, que es
precisamente la propuesta del trabajo: informe, evidencia visual y fiabilidad no deben consultarse por
separado. Y tercero, su idoneidad como \textbf{prototipo para un estudio de lectura} futuro, al
permitir presentar casos controlados y recoger la interacción del clínico.

\subsection{Un gradiente de confianza, no una alerta de error}
La decisión de diseño más relevante es presentar la fiabilidad como un \textbf{semáforo en gradiente}
—una banda continua de verde a rojo— en lugar de una \textbf{alerta binaria} de posible error. La
justificación es doble. Por un lado, un clasificador binario sería \emph{poco preciso} dada la baja
prevalencia de las alucinaciones: generaría muchas falsas alarmas por cada acierto (véase
Tabla~\ref{tab:oper}). Por otro, y más importante, la literatura de sistemas de ayuda a la decisión
clínica documenta el fenómeno de la \textbf{fatiga de alarmas}: cuando un sistema emite alertas
frecuentes —muchas de ellas irrelevantes— los profesionales se \emph{desensibilizan} y terminan
ignorándolas, anulando su utilidad e incluso introduciendo riesgo. Un gradiente calibrado evita este
efecto: la mayoría de los informes fiables permanecen en verde, de baja fricción, y la atención del
radiólogo se \textbf{prioriza} hacia la minoría de casos de riesgo elevado, en lugar de dispersarse en
una sucesión de alertas. La confianza se convierte así en un \emph{ordenador de la carga de revisión},
coherente con los puntos de operación de la Tabla~\ref{tab:oper}. La Figura~\ref{fig:semaforo}
representa el semáforo con las bandas calibradas sobre Open-I.

\begin{figure}[t]\centering
\includegraphics[width=0.9\linewidth]{fig_semaforo.png}
\caption{Semáforo de confianza en gradiente. Las bandas, definidas sobre los cuartiles de la
distribución de los informes correctos, se anotan con la probabilidad de alucinación observada. El
verde concentra los informes fiables (baja fricción) y el rojo prioriza la revisión.}
\label{fig:semaforo}
\end{figure}

La Figura~\ref{fig:demo} ilustra la salida de la demostración sobre dos ejemplos reales de MAIRA-2:
un informe con una alucinación —una fabricación de infiltrado— que el sistema sitúa en zona roja
(baja confianza), y un informe normal correcto que sitúa en zona verde.

\begin{figure}[t]\centering
\includegraphics[width=0.47\linewidth]{fig_demo_halluc.png}\hfill
\includegraphics[width=0.47\linewidth]{fig_demo_correct.png}
\caption{Salida de la demostración (MAIRA-2) con el anclaje por cajas. Izquierda: informe con una
alucinación de infiltrado, con baja confianza (zona roja). Derecha: informe normal correcto, con alta
confianza (zona verde).}
\label{fig:demo}
\end{figure}

\subsection{El radiólogo como decisor y la mitigación del sesgo de automatización}
El sistema se concibe explícitamente como \textbf{asistivo, no autónomo}: no emite un diagnóstico
final ni bloquea informes, sino que aporta evidencia (anclaje) y una estimación de fiabilidad para que
el \textbf{radiólogo decida}. Este diseño mitiga el \textbf{sesgo de automatización} —la tendencia a
aceptar acríticamente la salida de un sistema automático— de dos formas. Primera, el anclaje se
presenta como puntero de \emph{dónde} mirar, no como prueba de acierto, evitando que una explicación
plausible induzca confianza injustificada. Segunda, la leyenda del semáforo incorpora una
\textbf{advertencia honesta}: el color mide la duda del modelo, de modo que el verde indica baja
probabilidad de \emph{alucinación} pero \emph{no} garantiza la ausencia de \emph{omisiones}, que el
modelo comete con confianza. Comunicar esta limitación en la propia interfaz es una salvaguarda de
seguridad, no un detalle menor. Finalmente, la calibración es \textbf{específica de cada modelo}
—cada generador dispone de su propia calibración del semáforo—, lo que reconoce que la relación entre
confianza y acierto no es universal y debe establecerse por sistema antes de su despliegue.

\section{Consideraciones éticas, legales y de sostenibilidad}
El despliegue de un sistema como el propuesto en un entorno sanitario real está sujeto a un marco
ético y normativo exigente que conviene explicitar. Desde el punto de vista \textbf{legal y
regulatorio}, un sistema que asiste al diagnóstico se considera \emph{producto sanitario} y queda
sujeto a la regulación europea correspondiente (Reglamento de Productos Sanitarios) y, de forma
creciente, al Reglamento de Inteligencia Artificial de la Unión Europea, que clasifica los sistemas de
apoyo a la decisión clínica como de \textbf{alto riesgo} y les impone requisitos de transparencia,
supervisión humana, gestión de riesgos y trazabilidad. El diseño asistivo del presente trabajo —con el
radiólogo como decisor y una salida de confianza explicable y calibrada— se alinea con estos
principios, especialmente con el de \textbf{supervisión humana significativa}. En cuanto a la
\textbf{protección de datos}, el trabajo emplea exclusivamente conjuntos de acceso libre y
anonimizados, evitando el tratamiento de datos personales; un despliegue clínico requeriría, además,
las garantías del Reglamento General de Protección de Datos.

Desde el punto de vista \textbf{ético}, dos riesgos merecen atención. El primero es el \textbf{sesgo}:
un modelo entrenado sobre poblaciones concretas puede rendir de forma desigual en otras, lo que este
trabajo aborda parcialmente al evaluar la generalización a un dominio distinto y al recomendar una
calibración específica por sistema y población antes del despliegue. El segundo es el ya mencionado
\textbf{sesgo de automatización}, cuya mitigación se integra en el propio diseño de la interfaz. En
cuanto a la \textbf{sostenibilidad}, el enfoque prioriza la \textbf{eficiencia}: todos los modelos se
ejecutan cuantizados en una única tarjeta gráfica de consumo, y la señal de fiabilidad propuesta se
obtiene \emph{sin coste adicional de inferencia} —reutiliza las probabilidades que el modelo ya
produce—, evitando el gasto computacional de métodos que exigen múltiples generaciones o modelos
auxiliares. Esta frugalidad favorece la accesibilidad del sistema en entornos con recursos limitados.

\section{Impacto, agentes implicados y viabilidad}
Los \textbf{agentes} implicados en un despliegue del sistema son el \textbf{radiólogo} —usuario y
decisor—, el \textbf{servicio de radiología} —que gestiona la carga de trabajo—, el \textbf{paciente}
—beneficiario último de una lectura más segura— y el \textbf{equipo técnico} responsable de la
calibración y el mantenimiento. El \textbf{impacto} potencial es triple. En el plano
\textbf{asistencial}, un priorizador de revisión bien calibrado permite concentrar la atención del
radiólogo en los informes de mayor riesgo, mitigando la presión asistencial sin comprometer la
seguridad. En el plano de la \textbf{seguridad del paciente}, la comunicación explícita y honesta de
la fiabilidad —incluida la advertencia sobre las omisiones— reduce el riesgo de una confianza
injustificada en la salida automática. Y en el plano \textbf{científico-técnico}, el trabajo aporta
una metodología de auditoría de la fiabilidad reutilizable para otros modelos y dominios.

En cuanto a la \textbf{viabilidad}, el enfoque es deliberadamente \textbf{frugal}: todos los modelos
se ejecutan cuantizados en una única tarjeta gráfica de consumo, sin infraestructura especializada, y
la señal de confianza no incurre en coste de inferencia adicional. El principal recurso consumido es
el \textbf{tiempo de cómputo} de la extracción de señales y el etiquetado, mitigado mediante procesos
reanudables. Esta frugalidad hace el sistema económicamente viable y accesible en entornos con
recursos limitados, y compatible con su integración como módulo asistivo en flujos de trabajo
existentes.

\section{Discusión y análisis de los resultados}
Los resultados dibujan una imagen coherente. La explicación visual de estos modelos es \emph{fiel y
está anclada}, lo que valida su uso como puntero de ``dónde mirar''; pero la fiabilidad \textbf{no se
lee en la forma de la explicación} —ni en el mapa de calor, ni en el \emph{lift} objetivo, ni en las
cajas entrenadas del propio modelo— sino en la \textbf{incertidumbre de la decisión}, y de forma
robusta solo para las alucinaciones. Esta señal, una vez depurada del token de apertura, resulta
\emph{simple, agnóstica al modelo y disponible en inferencia}, y —a diferencia de lo que sugería el
único generador inicial— \emph{se transfiere fuera de distribución} cuando el generador es de calidad
suficiente. Que un modelo mejor mantenga la señal fuera de distribución (0,81) mientras el débil la
pierde (0,61) indica que el problema no es el método sino el generador que, sin datos del dominio,
alucina con confianza. La exploración de señales adicionales (entropía, verosimilitud,
auto-consistencia, dependencia de imagen) mostró que ninguna mejora de forma consistente en los tres
datasets: son específicas de modelo o de dominio, por lo que la fórmula final privilegia la
\emph{simplicidad robusta} de una única señal frente a un combinado sobreajustado.

Las omisiones quedan como problema de \textbf{sensibilidad} del generador, detectable solo en parte
desde sus propias señales y por otra vía (mejor generador o anclaje explícito). Como aportación
transversal, de carácter metodológico, conviene subrayar que repetidamente señales de apariencia
fuerte (el máximo de atención, la entropía cruda confundida por la longitud, la relación
\emph{lift}--acierto, las propias cajas del modelo) resultaron artefactos que solo un control
cuidadoso de confusores y la validación contra referencias objetivas revelaron; en este dominio,
\emph{validar contra verdad objetiva es imprescindible}.

En relación con la literatura, estos resultados matizan y extienden el conocimiento previo. Confirman,
por un lado, la intuición de que la incertidumbre del modelo es un mejor predictor de la alucinación
que sus explicaciones visuales, en línea con los métodos de auto-consistencia y de estimación de
incertidumbre desarrollados para modelos de lenguaje generales; pero aportan, por otro, tres elementos
propios del dominio radiológico que la literatura no había establecido de forma conjunta: la
\textbf{disociación} cuantitativa entre alucinaciones detectables y omisiones que no lo son, la
\textbf{robustez cruzada} de la señal entre modelos y dominios —con la identificación del generador
débil como causa de la degradación fuera de distribución— y la constatación, con las cajas de anclaje
del propio modelo, de que ni siquiera una explicación entrenada y plausible predice el acierto. El
valor práctico de estos hallazgos es directo: indican que un sistema de ayuda al radiólogo debe
construir su señal de confianza sobre la incertidumbre de la decisión, comunicar honestamente su
alcance (que no cubre las omisiones) y calibrarse por modelo antes de su despliegue.

Finalmente, procede señalar las limitaciones:
MAIRA-2 fue entrenado con PadChest, por lo que su evaluación in-distribution puede incluir solape
train--test —mitigado usando Open-I, ajeno a su entrenamiento, como prueba fuera de distribución—; el
número de alucinaciones por celda es reducido, lo que ensancha los intervalos de confianza; y el
beneficio clínico real de presentar informe, anclaje y confianza al radiólogo requeriría un estudio
prospectivo de lectura, que se propone como trabajo futuro.
